\section{Modality-Aware Partitioner}
\label{section:modality-aware-partitioning}

Building on the principles in \S\ref{section:design-decisions},
we propose \emph{modality-aware partitioning}.
This approach partitions each modality module into multiple model chunks before training,
and dynamically splits microbatches into modality-specific sub-microbatches during training.
Specifically, for the $i$-th modality, {\sys} determines two parameters:
the sub-microbatch size $B_i$ and the number of pipeline segments $K_i$.
Consequently, the $i$-th modality module is partitioned into $P \cdot K_i$ model chunks,
distributed across $P$ pipeline ranks.

\heading{Determine Sub-Microbatch Size}
\sys{} independently selects the minimum viable sub-microbatch size $B_i$ for each modality module through systematic profiling.
Although smaller sub-microbatch sizes enable finer-grained pipeline partitioning and higher scheduling efficiency,
excessively small sizes may lead to GPU underutilization (\figref{figure:pipeline-scheduling}).
To balance these trade-offs, \sys{} measures the computation latency of each modality module across varying sub-microbatch sizes.
The optimal $B_i$ is determined as the smallest sub-microbatch size that maintains at least 95\%
of the peak GPU computational efficiency observed across all tested configurations.

\heading{Partition Model Chunks}
After determining sub-microbatch sizes, \sys{} partitions each modality module to achieve global latency balancing.
Given $n$ modality modules with sorted computation latencies $T_1 \leq T_2 \leq \cdots \leq T_n$
(measured at their respective $B_i$ sizes),
\sys{} assigns pipeline segment counts proportional to these latencies.
Specifically, the fastest module ($T_1$) is configured as a single pipeline segment,
while the $i$-th module is partitioned into $K_i = \lfloor T_i / T_1 \rfloor$ segments.
For a modality module comprising $L_i$ layers,
\sys{} distributes layers across $P \cdot K_i$ model chunks,
with each chunk containing $L_i/(P \cdot K_i)$ consecutive layers.

\heading{Construct Sub-Microbatch}
During training execution, \sys{} dynamically constructs sub-microbatches for each modality module
based on the content of incoming microbatches.
For a microbatch assigned to the $i$-th modality module with $N_i$ instances,
\sys{} splits it into $M_i = \lceil N_i/B_i \rceil$ uniformly partitioned sub-microbatches.
This process generates $2 M_i \cdot K_i$ pipeline segments in total
(accounting for both forward and backward computations) for the modality module.

\section{Pipeline Schedule Searcher}
\label{section:pipeline-schedule-searcher}

\subsection{MCTS-Based Pipeline Segment Reordering}
\label{section:segment-reordering}

\sys{} determines the relative ordering between pipeline segments by assigning scheduling priorities to them.
These priorities are subsequently utilized in the pipeline stage interleaving phase (\S\ref{section:stage-interleaving}).
To efficiently explore the search space of possible priority assignments, \sys{} employs Monte Carlo tree search (MCTS)~\cite{mcts}.

\heading{Search Tree Construction}
Given $n$ pipeline segments, MCTS constructs a sequence of length $n$ by iteratively selecting the segment for each position $i$ (from $1$ to $n$).
The segment at the $i$-th position receives a priority of $n - i$, ensuring earlier selected segments obtain higher priorities.

The algorithm builds a search tree where each node at depth $d$ corresponds to the $d$-th element in the sequence.
Consequently, any root-to-node path ending at depth $d$ represents a partial sequence for the first $d$ positions.
Each node $v$ maintains the highest performance score $s_v$ observed among its descendants in the MCTS tree.

\heading{Search Loop}
MCTS iteratively performs the following four phases until convergence:

\begin{enumerate}
    \item \textbf{Node Selection:} Starting from the root, MCTS navigates downward by selecting child nodes that maximize the upper confidence bound (UCB) score $s_v^\alpha + \beta\sqrt{(\log{N_x}) / N_v}$.
    Here, $x$ denotes the current node, $v$ is the target child node, $N_x$ and $N_v$ are visit counts, and $\alpha$, $\beta$ are hyperparameters.
    This process repeats until a leaf node $u$ is reached.

    \item \textbf{Tree Expansion:} A new node representing the next element after $u$ is added to the search tree.
    Then $u$ is set to the newly created node.

    \item \textbf{Random Rollouts:} Multiple rollouts (\eg 10 trials) generate complete sequences by randomly assigning segments to remaining positions after $u$.
    Each sequence then undergoes pipeline stage interleaving (\S\ref{section:stage-interleaving}) and per-layer memory optimization (\S\ref{section:memory-optimization}) to compute performance scores (\ie end-to-end iteration time).

    \item \textbf{Score Backpropagation:} The best rollout score among all trials propagates upward from $u$ to update the performance scores of ancestor nodes.
\end{enumerate}

This loop continues until either a predefined time budget is exhausted or the entire search space is explored.

\heading{Optimization}
We observe that segments processing the same modality within a single microbatch have identical pipeline structures and similar stage latencies.
Consequently, their relative ordering does not impact end-to-end performance.
This insight allows \sys{} to assign identical priorities to such segments and enforce a fixed ordering between them, significantly reducing the search space.

\subsection{Greedy Pipeline Stage Interleaving}
\label{section:stage-interleaving}

After assigning priorities to segments,
\sys{} employs a dual-queue algorithm to adaptively interleave forward and backward pipeline stages.
When both schedulable forward and backward stages are available,
\sys{} mimics Megatron-LM's memory-efficient ``one-\-forward-\-one-\-backward'' (1F1B) pattern~\cite{megatron-lm-arxiv20,dynapipe-eurosys24}.
When either type is unavailable,
a greedy strategy fills pipeline bubbles to construct compact schedules.

\heading{Initialization}
\sys{} first obtains the stage latencies and memory consumptions for all stages via the training simulator,
using the most memory-efficient scheme determined in \S\ref{section:memory-optimization}.
This step ensures sufficient optimization space for subsequent per-layer memory optimizations.

For each pipeline rank, \sys{} maintains:
(1) the end time of the last scheduled stage ($t_\mathrm{last}$, initially zero);
(2) two priority queues ($Q_\mathrm{fw}$, $Q_\mathrm{bw}$) storing forward and backward stages in descending priority order; and
(3) the minimum start time ($t_\mathrm{min}$) among stages in $Q_\mathrm{fw}$ and $Q_\mathrm{bw}$.
Each stage maintains a minimum schedulable start time ($t_\mathrm{start}$) initialized to zero for stages with no predecessors and $+\infty$ otherwise.

\heading{Iterative Scheduling}
\sys{} iteratively selects a pipeline rank and schedules one stage from it:
\begin{enumerate}
    \item Select the pipeline rank with the smallest $t_\mathrm{min}$.
    \item Compare the minimum start times of stages in $Q_\mathrm{fw}$ ($t_\mathrm{fw}$) and $Q_\mathrm{bw}$ ($t_\mathrm{bw}$) against $t_\mathrm{last}$.
    \item If both $t_\mathrm{fw} < t_\mathrm{last}$ and $t_\mathrm{bw} < t_\mathrm{last}$,
          schedule alternating forward/backward stages based on the last scheduled stage's computation type,
          emulating the 1F1B pattern.
    \item Otherwise, select the stage with the smallest $t_\mathrm{start}$ to minimize pipeline bubble between $t_\mathrm{min}$ and stage's start time.
\end{enumerate}

The selected stage is dequeued and scheduled on the pipeline rank.
Subsequently, $t_\mathrm{last}$, $t_\mathrm{min}$, and $t_\mathrm{start}$ values for all successor stages are updated.

\heading{Memory Constraints}
Throughout scheduling, \sys{} tracks real-time memory consumption across pipeline ranks.
When a rank exceeds memory capacity, its forward queue is temporarily disabled to prevent memory overflow.

\subsection{Per-Layer Memory Optimization}
\label{section:memory-optimization}

\sys{} adaptively selects appropriate memory optimization strategies (\eg activation checkpointing)
for all model layers.
Since the stage interleaving scheme is predetermined by the dual-queue algorithm (\S\ref{section:stage-interleaving}),
this phase independently optimizes end-to-end latency for each pipeline rank.

\heading{Offline Candidate Generation}
Before training, \sys{} enumerates all possible optimization strategies for each model layer
and estimates their execution time and memory consumption using a training simulator.
The system then groups each forward stage with its corresponding backward stage into a \emph{stage pair}.
For each stage pair, \sys{} selects up to $S$ candidate strategies (\eg $S = 10$)
from the combinatorial space of layer-wise strategies through a three-step process:
(1) identifying the fastest candidate,
(2) identifying the most memory-efficient candidate, and
(3) evenly partitioning the memory range between these extremes into $S-2$ buckets,
then selecting the most time-efficient candidate within each bucket via a multiple-choice knapsack algorithm~\cite{multi-choice-knapsack}.

\heading{ILP Formulation}
For each pipeline rank, \sys{} employs an ILP solver to select optimal candidates.
Given $n$ stage pairs, each with $S$ candidate strategies, let $s_i$, $t_i$ denote the start and end timestamps of the $i$-th stage pair,
and $\mathrm{lat}_{i,j}$, $\mathrm{mem}_{i,j}$ denote the stage latency and memory consumption of the $j$-th candidate for the $i$-th stage pair.
The formulation uses:

\begin{itemize}
    \item \textbf{Variables:} $o_{i,j} \in \{0,1\}$ indicating whether the $j$-th candidate is selected for the $i$-th stage pair.
    \item \textbf{Objective:} Minimize total latency $\sum_i^n \sum_j^S o_{i,j} \cdot \mathrm{lat}_{i,j}$.
    \item \textbf{Selection Constraints:} Each stage pair selects exactly one strategy, \ie $\sum_j^S o_{i,j} = 1$ ($\forall 1 \leq i \leq n$).
    \item \textbf{Memory Constraints: } At any time $s_k$, the memory limit $M$ is satisfied,
    \ie $\sum_i^n [s_i \leq s_k \leq t_i] \sum_j^S o_{i,j} \cdot \mathrm{mem}_{i,j} \leq M$ ($\forall 1 \leq k \leq n$),
    where $[\cdot]$ is the Iverson bracket~\cite{iverson-bracket}.
\end{itemize}

\heading{Optimizations}
Although the problem scale has been significantly reduced compared to the monolithic, end-to-end ILP formulation,
solving an ILP instance may still require several seconds.
We introduce two techniques to achieve efficient solving ($<$10 ms per instance on a single CPU core):
(1) warm-starting the ILP solver with a greedy initial solution, and
(2) permitting a small optimality gap (\eg $\leq 5\%$) for early termination,
as closing the final 5\% gap incurs diminishing returns but prohibitive computational costs.

\subsection{Time Complexity Analysis}

In this section, we outline the algorithmic structure of {\sys}'s training planner and
analyze its time complexity.
The planner operates as an iterative search loop that continues until a predefined time budget is exhausted.
Each iteration consists of three stages:
pipeline segment reordering (\S\ref{section:segment-reordering}),
pipeline stage interleaving (\S\ref{section:stage-interleaving}), and
per-layer memory optimization (\S\ref{section:memory-optimization}).
For the complexity analysis of a single search iteration, let
$p$ denote the number of pipeline ranks,
$n$ the number of sub-microbatches, and
$S$ the number of candidate memory-saving strategies.
Additionally, let $\mathrm{ApproxILP}(m, k)$ represent the time complexity of the approximate ILP solver for
an instance with $\mathrm O(m)$ variables and $\mathrm O(k)$ constraints.

First, pipeline segment reordering generates an ordering for all sub-microbatches in $\mathrm O(n)$ time.
Second, the greedy scheduling algorithm for pipeline stage interleaving runs in time proportional to the total number of pipeline stages, \ie $\mathrm O(p \cdot n)$.
Finally, per-layer memory optimization solves a set of small, approximate ILP instances to select the optimal configuration for each pipeline rank.
Each ILP instance requires creating $n \cdot S$ indicator variables ($o_{i,j}$),
along with $n$ selection constraints and $n$ memory constraints.
Consequently, the time complexity for per-layer memory optimization is $\mathrm O(p \cdot \mathrm{ApproxILP}(n \cdot S, n))$.
Combining all these steps, the overall time complexity for a single search iteration is $\mathrm O(p \cdot (n + \mathrm{ApproxILP}(n \cdot S, n)))$.

In contrast, a full ILP baseline approach is inherently inefficient due to the massive scale of variables and constraints involved.
First, it requires solving the entire pipeline globally, rather than decomposing the problem per pipeline rank.
Second, it necessitates $\mathrm O(n^2)$ ordering constraints to prevent overlapping stage executions within each pipeline rank, whereas {\sys} avoids this overhead through efficient explicit scheduling.
Furthermore, in the memory optimization phase, assuming $L$ layers per pipeline stage and $c$ candidate strategies per layer, the search space expands to $c^L$ candidates per stage.
Consequently, the complexity of the baseline ILP approach escalates to $\mathrm O(\mathrm{ApproxILP}(p \cdot n \cdot c^L, p \cdot n^2))$.
We empirically compare the performance of this baseline against {\sys}'s training planner in \S\ref{section:planner-evaluation} and \figref{figure:search-scalability}.
